% GENERATED FILE. Edit canonical lessons, then run npm run generate:books.
% canonical-source-hash: fnv1a64:b3a029ddf3c2fa8a
% canonical-lessons: FA-C30-ghalb, FA-C30-bazu, FA-C30-tab, FA-C30-bimarestan, FA-C30-bimar

\chapter{Heart, Arm, Fever, Hospital, Ill}
\label{ch:fa-heart-arm-fever-hospital-ill}
\hlchaptermodality{30}

\begin{chapteropening}
\textbf{By the end of this chapter:} \emph{I can say the heart, an arm, a fever, a hospital and ill, sick.}

Complete the last lesson of chapter 30: I can say the heart, an arm, a fever, a hospital and ill, sick.
\end{chapteropening}

\section[\texorpdfstring{\fa{قلب} (ghalb)}{قلب (ghalb)}]{\fa{قلب} (ghalb) — the heart}

\label{lesson:FA-C30-ghalb}



\begin{quote}
Before the new one: say the Persian for the belly, then the Persian for a lip.
\end{quote}

\subsection*{You'll want to know: \fa{قلب}}
\textbf{\fa{قلب}} — \emph{ghalb} — \textquotedblleft{}the heart\textquotedblright{}.

\begin{cousinweb}
\emph{ghalb} is borrowed from Arabic \emph{qalb}, \textquotedblleft{}heart\textquotedblright{}. Its Arabic root is \emph{q-l-b} (Lane). The root is about turning: Lane gives \textquotedblleft{}he turned it upside-down\textquotedblright{} for the verb.
\end{cousinweb}

\begin{grammarlens}[title={What you've built}]
One more word for this chapter.
\end{grammarlens}

\subsection*{Your turn}
\begin{itemize}
\raggedright
  \item \emph{Say it:} \emph{ghalb}
  \item \emph{Say it:} \emph{ghalb}, once more
  \item \emph{Say it:} say \emph{ghalb}
  \item \emph{Recall:} say the Persian for the belly, then the Persian for a lip, then say \emph{ghalb} again
\end{itemize}

\begin{tcolorbox}[breakable,colback=teal!4,colframe=teal!35!black,title={Before you move on}]
What does \fa{قلب} mean? (\textquotedblleft{}The heart\textquotedblright{}.) Say it once more.
\end{tcolorbox}
\section[\texorpdfstring{\fa{بازو} (bâzu)}{بازو (bâzu)}]{\fa{بازو} (bâzu) — an arm}

\label{lesson:FA-C30-bazu}



\begin{quote}
Before the new one: say the Persian for a lip, then the Persian for the heart.
\end{quote}

\subsection*{You'll want to know: \fa{بازو}}
\textbf{\fa{بازو}} — \emph{bâzu} — \textquotedblleft{}an arm\textquotedblright{}.

\begin{cousinweb}
\emph{bâzu} is a native Persian word, inherited rather than borrowed; it is usually traced to Middle Persian \emph{bāzūg}, \textquotedblleft{}arm\textquotedblright{}. Old Persian had \emph{bāzu-}, \textquotedblleft{}arm\textquotedblright{} (EtymDB), and EtymDB takes it back to Indo-European \emph{*b\textsuperscript{h}eh\textsubscript{2}ǵ\textsuperscript{h}ús}, \textquotedblleft{}arm\textquotedblright{}; English \emph{bough}, which once meant \textquotedblleft{}shoulder\textquotedblright{}, is usually traced to the same word.
\end{cousinweb}

\begin{grammarlens}[title={What you've built}]
One more word for this chapter.
\end{grammarlens}

\subsection*{Your turn}
\begin{itemize}
\raggedright
  \item \emph{Say it:} \emph{bâzu}
  \item \emph{Say it:} \emph{bâzu}, once more
  \item \emph{Say it:} say \emph{bâzu}
  \item \emph{Recall:} say the Persian for a lip, then the Persian for the heart, then say \emph{bâzu} again
\end{itemize}

\begin{tcolorbox}[breakable,colback=teal!4,colframe=teal!35!black,title={Before you move on}]
What does \fa{بازو} mean? (\textquotedblleft{}An arm\textquotedblright{}.) Say it once more.
\end{tcolorbox}
\section[\texorpdfstring{\fa{تب} (tab)}{تب (tab)}]{\fa{تب} (tab) — a fever}

\label{lesson:FA-C30-tab}



\begin{quote}
Before the new one: say the Persian for the heart, then the Persian for an arm.
\end{quote}

\subsection*{You'll want to know: \fa{تب}}
\textbf{\fa{تب}} — \emph{tab} — \textquotedblleft{}a fever\textquotedblright{}.

\begin{cousinweb}
\emph{tab} is a native Persian word, inherited rather than borrowed; it is usually traced to Middle Persian \emph{tab}, \textquotedblleft{}fever\textquotedblright{}. It belongs with the old verb root \emph{tap-}, \textquotedblleft{}to be hot\textquotedblright{}.
\end{cousinweb}

\begin{grammarlens}[title={What you've built}]
One more word for this chapter.
\end{grammarlens}

\subsection*{Your turn}
\begin{itemize}
\raggedright
  \item \emph{Say it:} \emph{tab}
  \item \emph{Say it:} \emph{tab}, once more
  \item \emph{Say it:} say \emph{tab}
  \item \emph{Recall:} say the Persian for the heart, then the Persian for an arm, then say \emph{tab} again
\end{itemize}

\begin{tcolorbox}[breakable,colback=teal!4,colframe=teal!35!black,title={Before you move on}]
What does \fa{تب} mean? (\textquotedblleft{}A fever\textquotedblright{}.) Say it once more.
\end{tcolorbox}
\section[\texorpdfstring{\fa{بیمارستان} (bimârestân)}{بیمارستان (bimârestân)}]{\fa{بیمارستان} (bimârestân) — a hospital}

\label{lesson:FA-C30-bimarestan}



\begin{quote}
Before the new one: say the Persian for an arm, then the Persian for a fever.
\end{quote}

\subsection*{You'll want to know: \fa{بیمارستان}}
\textbf{\fa{بیمارستان}} — \emph{bimârestân} — \textquotedblleft{}a hospital\textquotedblright{}.

\begin{cousinweb}
\emph{bimârestân} is built from \emph{bimâr}, \textquotedblleft{}sick\textquotedblright{}, and the place suffix \emph{-estân}: \textquotedblleft{}the place of the sick\textquotedblright{}. Both pieces are inherited: \emph{bimâr} continues Middle Persian \emph{wēmār} (EtymDB), and \emph{-estân} is the Persian suffix of place names that \emph{Pakistan} also uses.
\end{cousinweb}

\begin{grammarlens}[title={What you've built}]
One more word for this chapter.
\end{grammarlens}

\subsection*{Your turn}
\begin{itemize}
\raggedright
  \item \emph{Say it:} \emph{bimârestân}
  \item \emph{Say it:} \emph{bimârestân}, once more
  \item \emph{Say it:} say \emph{bimârestân}
  \item \emph{Recall:} say the Persian for an arm, then the Persian for a fever, then say \emph{bimârestân} again
\end{itemize}

\begin{tcolorbox}[breakable,colback=teal!4,colframe=teal!35!black,title={Before you move on}]
What does \fa{بیمارستان} mean? (\textquotedblleft{}A hospital\textquotedblright{}.) Say it once more.
\end{tcolorbox}
\section[\texorpdfstring{\fa{بیمار} (bimâr)}{بیمار (bimâr)}]{\fa{بیمار} (bimâr) — ill, sick}

\label{lesson:FA-C30-bimar}



\begin{quote}
Before the new one: say the Persian for a fever, then the Persian for a hospital.
\end{quote}

\subsection*{You'll want to know: \fa{بیمار}}
\textbf{\fa{بیمار}} — \emph{bimâr} — \textquotedblleft{}ill, sick\textquotedblright{}.

\begin{cousinweb}
\emph{bimâr} is a native Persian word, inherited rather than borrowed: it continues Middle Persian \emph{wēmār}, \textquotedblleft{}sick, ill\textquotedblright{} (EtymDB).
\end{cousinweb}

\begin{grammarlens}[title={What you've built}]
One more word for this chapter.
\end{grammarlens}

\subsection*{Your turn}
\begin{itemize}
\raggedright
  \item \emph{Say it:} \emph{bimâr}
  \item \emph{Say it:} \emph{bimâr}, once more
  \item \emph{Say it:} say \emph{bimâr}
  \item \emph{Recall:} say the Persian for a fever, then the Persian for a hospital, then say \emph{bimâr} again
\end{itemize}

\begin{tcolorbox}[breakable,colback=teal!4,colframe=teal!35!black,title={Before you move on}]
What does \fa{بیمار} mean? (\textquotedblleft{}Ill, sick\textquotedblright{}.) Say it once more.
\end{tcolorbox}
